\subsection{Propósito}

Este documento define el estándar de codificación en C\# para el desarrollo de un videojuego de mesa en dos dimensiones. Su propósito es establecer un conjunto único de reglas de nombrado, estilo, complejidad, manejo de errores, documentación, bitácoras y pruebas unitarias que todo el equipo debe seguir, de manera que el código resultante sea consistente, legible y mantenible independientemente de quién lo haya escrito.

El estándar es deliberadamente independiente de cualquier motor de videojuegos (Unity, Godot, Unreal u otro) y de cualquier tecnología de interfaz gráfica o de conexión de red: las reglas aquí descritas aplican al lenguaje C\# en sí mismo y a las bibliotecas oficiales de .NET, por lo que son válidas sin importar la tecnología concreta sobre la que finalmente se implemente el juego. Cada regla incluye, además, la justificación técnica o académica en la que se apoya, para que el equipo comprenda no solo \emph{qué} debe hacer, sino \emph{por qué} conviene hacerlo así.
